\section{One shear and exact coverage}\label{sec:final}

We extend the simultaneous Boolean-shift argument of Li and Liu
\cite[Theorem~4.3]{LL} to patterns whose branch bases may depend on later
coordinates. Unit separation between siblings still allows one linear
shear to treat every pattern in a finite family. Throughout this section
all measures on a lattice torus are normalized Haar measures.

\subsection{A simultaneous shear for finitely many patterns}

A labeled pattern in $\R^q$ is a finite nonempty set of points, with one
label assigned to each point. Call it \emph{binary by suffixes} if,
after fixing any suffix of coordinates, the next coordinate has at most
two possible values, and two such values always differ by $1$.
For $q=0$ a pattern is a single label. Such a pattern has at most $2^q$
points. Deleting points preserves the suffix rule; the result is again a
pattern whenever it is nonempty. Taking a nonempty slice at the last
coordinate and then deleting that coordinate also preserves the rule.

Figure~\ref{fig:suffix-tree} shows why a fixed global binary cube is not
required.

\begin{lemma}[Simultaneous shear]\label{lem:shear}
Let $\mathcal T=\R^m/\Lambda_h$, and let $U_a\subset\mathcal T$ be
measurable sets with $|U_a|\le\eta_a$, where $\eta_a>0$.
For $0\le q\le D$, let $\mathcal P_q$ be finite families of labeled
patterns in $\R^q$, binary by suffixes. Suppose they are closed under
last-coordinate slicing. Write
\[
 Q=\sum_{q=0}^D|\mathcal P_q|,\qquad L_*>\max(1,Q).
\]
There are real vectors $Z_1,\ldots,Z_D\in\R^m$ such that, for every
$P\in\mathcal P_q$,
\begin{equation}\label{eq:shear-bound}
 \left|\bigcap_{l\in P}
   \left(U_{a(l)}+
     \pi_{\Lambda_h}\left(\sum_{j=1}^q l[j]Z_j\right)\right)\right|
 \le L_*^{|P|-1}\prod_{l\in P}\eta_{a(l)}.
\end{equation}
Here $l[j]$ denotes the $j$th scalar coordinate of $l$.
\end{lemma}

\begin{proof}
The assertion in dimension zero is the assumed bound on each $U_a$.
Suppose $Z_1,\ldots,Z_{q-1}$ have been fixed so that all preceding
estimates hold. A pattern with just one value of $l[q]$ preserves its
bound under the resulting common translation.

Otherwise the two last-coordinate values are $v$ and $v+1$.
Let the two slices have cardinalities $k_1$ and $k_2$, and let $A$ and
$B$ be their intersection sets. Choose $Z_q$ uniformly in a fixed
fundamental parallelepiped of $\Lambda_h$. Translation invariance gives
\[
 |(A+\pi(vZ_q))\cap(B+\pi((v+1)Z_q))|
       =|A\cap(B+\pi(Z_q))|.
\]
This identity is valid even when $v$ is not an integer: the products
$vZ_q$ and $(v+1)Z_q$ are formed in $\R^m$ before projection.
Only the relative shift must be Haar uniform. By Fubini,
\[
 \E_{Z_q}|A\cap(B+\pi(Z_q))|=|A||B|
 \le L_*^{k_1+k_2-2}\prod_{l\in P}\eta_{a(l)}.
\]
Markov's inequality bounds the probability of violating
Equation~\eqref{eq:shear-bound} by $1/L_*$. If either slice has
measure zero the intersection has measure zero almost surely,
which gives the same conclusion. A union bound over
$\mathcal P_q$ has total probability less than one, so one choice of
$Z_q$ works for all its patterns. Induction proves the result.
\end{proof}

\begin{figure}[htbp]
\centering
\begin{tikzpicture}[x=1cm,y=1cm,
  every node/.style={font=\small},
  edge/.style={-{Latex[length=1.6mm]},thin}]
\node (root) at (0,0) {$P$};
\node (left) at (-2.7,-1.2) {$P_v$};
\node (right) at (2.7,-1.2) {$P_{v+1}$};
\draw[edge] (root) -- node[above left] {$l[q]=v$} (left);
\draw[edge] (root) -- node[above right] {$l[q]=v+1$} (right);
\node (la) at (-4,-2.5) {$P_{v,a}$};
\node (lb) at (-1.4,-2.5) {$P_{v,a+1}$};
\node (ra) at (1.4,-2.5) {$P_{v+1,b}$};
\node (rb) at (4,-2.5) {$P_{v+1,b+1}$};
\draw[edge] (left) -- node[above left] {$l[q-1]=a$} (la);
\draw[edge] (left) -- node[above right] {$a+1$} (lb);
\draw[edge] (right) -- node[above left] {$b$} (ra);
\draw[edge] (right) -- node[above right] {$b+1$} (rb);
\foreach \name in {la,lb,ra,rb}
  \draw[densely dotted] (\name.south) -- ++(0,-0.45);
\node at (0,-3.4) {Continue through the earlier coordinates.};
\end{tikzpicture}

\caption{The suffix rule used by the shear induction. The bases $a$ and
$b$ may differ because the later coordinate has already been fixed.
Only the separation of siblings is prescribed. Removing leaves, or an
entire subtree, preserves the rule.}
\label{fig:suffix-tree}
\end{figure}

\subsection{Removing a small hole}

The next argument is the classical completion-by-dilation method of
Rogers~\cite{Rog59}. We use the small-dilation form also given by
Fejes T\'oth~\cite[Lemma~2]{FT09}, and include its proof for the
single-lattice setting.

\begin{lemma}[Completion]\label{lem:completion}
Let $K\subset\R^n$ be a convex body and $\Gamma$ a full-rank lattice.
If its hole proportion $\delta=1-|\pi_\Gamma(K)|$ satisfies
\[
 \delta<(1-\delta)k^{-n}
\]
for a positive integer $k$, then $(1+1/k)K+\Gamma=\R^n$.
\end{lemma}

\begin{proof}
Multiplication by $k$ on $\R^n/\Gamma$ maps $\pi_\Gamma(K/k)$
onto $\pi_\Gamma(K)$. Its image has measure at most $k^n$ times
the measure of the original set: partition a fundamental
parallelepiped into $k^n$ cells on which the map is injective.
Consequently
\[
 |\pi_\Gamma(K/k)|\ge k^{-n}(1-\delta)>\delta.
\]
For every point $x$ of the torus, the sets $\pi_\Gamma(K)$ and
$x-\pi_\Gamma(K/k)$ have measures summing to more than one and
therefore intersect. Thus
$\pi_\Gamma(K)+\pi_\Gamma(K/k)$ is the whole torus.
Finally, convexity gives $K+K/k=(1+1/k)K$. Neither symmetry nor
the inclusion $0\in K$ is required.
\end{proof}

\subsection{Proof of the covering theorem}

\begin{proof}[Proof of Theorem~\ref{thm:main}]
All constants in the vertical construction, including its cutoff $B$,
are fixed first. For sufficiently large $n$, put $s=\log\log n$ and
use Proposition~\ref{prop:vertical} with $b_1=\lceil s^2\rceil$.
Retain its notation $L,M,D,D_v,b_i,h_i$, and set
\[
 R^2=\sum_{i=1}^{r+1}b_i h_i^2,\qquad m=n-D.
\]
Here $D=O(s^2+1)$ and $R$ is bounded by a fixed polynomial in $s$.
After passing to the affine coordinates of Lemma~\ref{lem:sections},
there are at most $(Cm^2)^D$ horizontal bodies $J_a$ such that every
$v$ with $\|v\|\le R$ has a label satisfying
\begin{equation}\label{eq:final-sections}
 J_a\subset K_v,\qquad
 c_{\rm sec}\vol_n(K)\gamma_D(v)
 \le\vol_m(J_a)\le
 C_{\rm sec}\vol_n(K)\gamma_D(v).
\end{equation}

Choose an absolute constant $C_*$ later and put
\begin{equation}\label{eq:target-density}
 \rho=C_*n\log n,\qquad
 D_h=\frac{\vol_n(K)}{\rho D_v},\qquad
 u_a=\frac{\vol_m(J_a)}{D_h}.
\end{equation}
For any vertical target $y$, use its pattern from
Proposition~\ref{prop:vertical} and label each raw point $l$ by a
body contained in the section at $y-Ml$. The residual bound permits
Equation~\eqref{eq:final-sections}. With the chosen horizontal covolume,
the normalized section volume is comparable to the corresponding
normalized Gaussian weight multiplied by $\rho$:
\[
 c_{\rm sec}\rho D_v\gamma_D(y-Ml)
 \le u_{a(l)}\le
 C_{\rm sec}\rho D_v\gamma_D(y-Ml).
\]
Thus the total weight supplies the sum of the horizontal loads, while
the atom bound controls each load separately. These estimates give
\begin{equation}\label{eq:final-loads}
 \sum_l u_{a(l)}\ge c_3\rho,\qquad
 \max_l u_{a(l)}\le U_n:=C\rho\exp(-c b_1^{0.70}),
\end{equation}
where $c_3=c_{\rm sec}\kappa>0$ is absolute.

Three scale comparisons are needed. First,
\begin{equation}\label{eq:load-upper-check}
 \frac{U_n}{m}
 \le C'\exp\bigl(s-c's^{1.4}\bigr)\longrightarrow0.
\end{equation}
Second, every pattern has at most $2^D$ points, so deleting those with
$u_{a(l)}<\sqrt m$ loses at most
\begin{equation}\label{eq:load-discard-check}
 2^D\sqrt m=n^{1/2+o(1)}=o(\rho).
\end{equation}
Third, the logarithm of the number of labels is
$O(D\log m)=o(\sqrt m)$. These statements follow from
$D=O(s^2+1)$ and $\log n=e^s$; all their constants are independent
of $K$. In particular the retained pattern is nonempty and has load
at least $c_3\rho/2$ for large $n$.

Restrict the label list to $\sqrt m\le u_a\le U_n$.
Lemma~\ref{lem:horizontal} gives one lattice $\Lambda_h$ of covolume
$D_h$ for which the hole sets
\[
 U_a=(\R^m/\Lambda_h)\setminus\pi_{\Lambda_h}(J_a)
\]
satisfy $|U_a|\le e^{-u_a/2}$ simultaneously. We next select one
shear that works for every retained pattern, not a separate shear
for each target.

\paragraph{A finite family for all targets.}
Initially restrict to $y=Mz$ with $z\in[0,1)^D$. This is enough
because $L$ contains $\Z^D$; after the final lattice is defined,
integer raw translations will reduce any target to this region.
The finite-alphabet conclusion of Proposition~\ref{prop:vertical}
places each scalar raw coordinate in $[-Q_{\rm raw},Q_{\rm raw}]$, where
$Q_{\rm raw}=1+(D+2)^D$. In block $i\le r$ it is a multiple of $1/p_i$;
in the terminal block it is an integer. The logarithmic alphabet size
is therefore at most
\[
 O(D\log(D+2)+\max_i\log p_i)=O(s^{c_0})
\]
for some fixed absolute exponent $c_0$.

Order scalar coordinates by increasing block index, and take
$\mathcal P_q$ to be \emph{all} labeled patterns on the first $q$
coordinate alphabets that are binary by suffixes, for $0\le q\le D$.
This includes every retained full pattern: the base in a block depends
only on later blocks, and deleting points cannot introduce a third
value at any suffix. The families are closed under slicing.

Let $\mathcal A_j$ be the alphabet of scalar coordinate $j$, and let
$N_{\rm lab}$ be the number of retained section labels. A labeled point
consists of a coordinate tuple and one such label, so
$N=N_{\rm lab}\prod_{j=1}^D|\mathcal A_j|$ bounds their number in
each of these coordinate spaces. Consequently
\[
 \log(N+1)=O(D\log n+s^{c_1}),\qquad
 Q:=\sum_{q=0}^D|\mathcal P_q|
       \le (D+1)(N+1)^{2^D}.
\]
For example, the last bound follows by encoding a set of at most
$2^D$ points as a list of that length, padded with a dummy symbol.
Take $L_*=2(Q+1)$. Then
\begin{equation}\label{eq:pattern-cost}
 2^D\log L_*
 \le C4^D(D\log n+s^{c_1})+C2^D\log(D+2)
 =o(n\log n).
\end{equation}
Indeed $4^D$ times any fixed polynomial in $s$ is
$\exp(O(s^2))=o(n)$.

\paragraph{The covering lattice.}
Apply Lemma~\ref{lem:shear} with $\eta_a=e^{-u_a/2}$ and set
$Zl=\sum_{j=1}^D l[j]Z_j$. Define
\begin{equation}\label{eq:full-lattice}
 \Gamma=\{(h+Zl,Ml):h\in\Lambda_h,\ l\in L\}.
\end{equation}
It is a full-rank lattice, being the image of the product lattice
$\Lambda_h\times L$ under an invertible block triangular linear
map. Its covolume is exactly $D_hD_v$, so
$\vol_n(K)/\det\Gamma=\rho$.

For a fixed $y$ and a retained labeled pattern, each $J_{a(l)}$
lies in the section $K_{y-Ml}$. The uncovered horizontal points of
$K+\Gamma$, viewed in $\mathcal T=\R^m/\Lambda_h$, therefore lie in
\[
 \bigcap_l\bigl(U_{a(l)}+\pi_{\Lambda_h}(Zl)\bigr).
\]
For $y\in M[0,1)^D$, Lemma~\ref{lem:shear} and the retained load give
the uniform bound
\begin{equation}\label{eq:uniform-holes}
 \text{horizontal hole fraction at }y
 \le\exp\left(2^D\log L_*-\frac{c_3\rho}{4}\right).
\end{equation}
For other $y$, reduce $M^{-1}y$ modulo $\Z^D$ and translate by the
corresponding point $(Zk,Mk)\in\Gamma$. This only translates the
horizontal section on its torus, so the same bound holds.

Let $F_h,F_v$ be fundamental parallelepipeds of $\Lambda_h$ and
$ML$. Their product is a measurable fundamental domain for
$\Gamma$: first reduce the vertical coordinate using a point
$(Zl,Ml)$, and then reduce the horizontal coordinate using
$\Lambda_h$. The union $K+\Gamma$ is closed, since translates of
the compact set $K$ by a lattice are locally finite. Its complement
is therefore measurable, and Fubini integrates
Equation~\eqref{eq:uniform-holes} over $F_v$. No measurable selection
of auxiliary patterns is needed. Equations~\eqref{eq:target-density}
and~\eqref{eq:pattern-cost} yield
\[
 \delta:=1-|\pi_\Gamma(K)|
 \le\exp\left(o(n\log n)-\frac{c_3C_*}{4}n\log n\right)
 \le n^{-2n}
\]
once $C_*$ is a sufficiently large absolute constant and then $n$
is sufficiently large. Since $n^{-2n}<(1-n^{-2n})n^{-n}$ for
$n\ge2$, Lemma~\ref{lem:completion} with $k=n$ gives
$(1+1/n)K+\Gamma=\R^n$. Thus
\[
 \Lambda=(1+1/n)^{-1}\Gamma
\]
is a single covering lattice for $K$, of density
\[
 \frac{\vol_n(K)}{\det\Lambda}
 =(1+1/n)^n\rho\le eC_*n\log n.
\]

\paragraph{Affine invariance and the remaining dimensions.}
Translations of a body do not change the covering property, and
invertible linear maps applied to both body and lattice preserve
density. Hence the affine normalization can be undone.

There remains only a fixed finite set of dimensions. In each such
dimension choose a maximum-volume simplex with vertices in $K$,
which exists by compactness and has positive volume. Put its vertices
at $0,e_1,\ldots,e_n$ by an affine map. Replacing the $j$th nonzero
vertex by any $x\in K$ changes the determinant to $x_j$, so maximality
gives $|x_j|\le1$. Thus
\[
 [0,1/n]^n\subset K\subset[-1,1]^n.
\]
The lattice $n^{-1}\Z^n$ covers by translates of the smaller cube
and therefore of $K$, with density at most $(2n)^n$.
The maximum of $(2n)^n/(n\log n)$ over the remaining finite set
can be absorbed in one absolute constant. This proves the theorem
for every $n\ge2$.
\end{proof}
